\section{Chern-connection regularization}
\label{sec:psh-chern-regularization}\index{Regularization!Chern connection}

Fix $C_0>0$ and $u\in\PSH(X,\omega)$ with $-C_0\leq u\leq0$. Distances, gradients and tensor norms use the Riemannian metric associated to $\omega$.

Let $R_tu$ be Demailly's regularization formed with the Chern connection of $\omega$. It averages $u$ over the tangent fiber with a fixed smooth nonnegative radial kernel of integral one. The averaging map is a smooth realization of the holomorphic part of the formal fiberwise expansion of the exponential map; it need not be genuinely holomorphic in the fiber. The following classical estimate controls the truncated transform, including its mixed base--fiber terms; its sources are recorded in \cref{sec:history-psh}.

\Needspace{23\baselineskip}
\begin{lemma}\label{lem:chern-kiselman-input}
There are constants $A,K>0$ and $t_*>0$, depending only on the fixed geometry and bound for $u$, such that $R_tu$ is smooth for $0<t<t_*$,
\[
 -C_0\leq R_tu\leq0,\qquad
 R_tu+Kt^2\downarrow u\quad\text{as }t\downarrow0.
\]
For $0<\delta<t_*$ and $c>0$, the bounded function
\[
 \Phi_{c,\delta}
   =\inf_{0<t\leq\delta}
       \left\{R_tu+K(t^2-\delta^2)+K(t-\delta)
                         -c\log(t/\delta)\right\}
\]
is quasi-plurisubharmonic and satisfies
\[
 \omega+\ddc\Phi_{c,\delta}\geq-(Ac+2K\delta)\omega.
\]
The infimum is attained at some $t>0$ at every point.
\end{lemma}

Only this classical curvature estimate is retained as an input. Its linear term $K(t-\delta)$ matters in the Hermitian case: it absorbs the mixed base--fiber terms in the complex Hessian of the regularization. The function $-c\log(t/\delta)$ tends to $+\infty$ as $t\downarrow0$, while $u$ is bounded, which explains attainment away from zero. We require neither smoothness of the minimizing radius nor smoothness of the transform across the region where that radius equals $\delta$.

We record two elementary estimates for the averaging operators so that their use does not require a K\"ahler volume identity.

\Needspace{11\baselineskip}
\begin{lemma}\label{lem:chern-smoothing-integral}
For all sufficiently small $t>0$,
\[
 \|R_tu-u\|_{L^1(\mathrm{d}V)}\leq Ct^2,
 \qquad \|\nabla R_tu\|_\infty\leq C/t.
\]
The constants are uniform for $-C_0\leq u\leq0$.
\end{lemma}
\begin{proof}
Let $z$ be local base coordinates and $\xi\in\mathbb C^n$ the fiber variable in a smooth orthonormal tangent frame. Write $\rho(|\xi|^2)$ for the fixed radial averaging kernel, normalized against Lebesgue measure $\mathrm{d}\xi$. In these coordinates the averaging map $F$ gives
\[
 R_tu(z)=\int u(F(z,t\xi))\rho(|\xi|^2)\,\mathrm{d}\xi,
 \qquad F(z,0)=z,
\]
where the kernel is supported in $|\xi|\leq1$. The map $F$ is smooth in these variables and its fiber differential at zero is an isomorphism. We first claim that the adjoint with respect to $\mathrm{d}V$ satisfies
\[
 R_t^*1=1+O(t^2)
\]
uniformly. To see this, insert a smooth partition of unity in the $z$ variable with each term compactly supported in a frame chart. For fixed $\xi$, change variables $x=F(z,t\xi)$ on a neighborhood of that support. The inverse map and its Jacobian have uniform Taylor expansions in $t$. The term of order zero in the transformed density is the chosen partition function at $x$. Its term of order one is linear in $\xi$. Integrating that term against the even radial kernel gives zero. The remaining density is bounded by $Ct^2\mathrm{d}V(x)$, since the second derivatives of all maps, partition functions and volume coefficients are bounded on these finitely many compact charts. Summing the partition functions proves the claim.

Consequently
\[
 \left|\int_X(R_tu-u)\,\mathrm{d}V\right|
  =\left|\int_Xu(R_t^*1-1)\,\mathrm{d}V\right|\leq Ct^2.
\]
The preceding lemma gives $R_tu-u\geq-Kt^2$. For a real function $a\geq-Kt^2$ one has $|a|\leq a+2Kt^2$. Apply this inequality and integrate to obtain the first estimate.

For the derivative bound, change variables in the fiber to express $R_tu(z)=\int_Xu(x)k_t(z,x)\,\mathrm{d}V(x)$. In local coordinates the smooth kernel is supported where $|x-z|\leq Ct$ and satisfies
\[
 |\nabla_z k_t(z,x)|\leq Ct^{-2n-1}.
\]
Indeed, the inverse fiber map and its Jacobian have uniformly bounded derivatives; differentiating the scaled radial argument introduces at most one factor $t^{-1}$ in addition to $t^{-2n}$. Smooth compact support of the kernel causes no boundary term. Integration over a ball of volume at most $Ct^{2n}$ gives $\int|\nabla_zk_t|\,\mathrm{d}V\leq C/t$. Multiplication by the uniform bound for $u$ proves the assertion.
\end{proof}

The first estimate uses cancellation of the odd first-order term of an averaging kernel. It does not use conservation of Hermitian Monge--Amp\`ere mass.


\subsection{The general curvature modulus}
\label{subsec:general-curvature-modulus}

Let $\theta$ now be any smooth real $(1,1)$-form and let $u\in\PSH(X,\theta)\cap L^\infty(X)$. Use the same averaging operators and metric.

\Needspace{10\baselineskip}
We first specify the classical input for this background. Apply the classical Chern-connection regularization theorem to the closed current $\ddc u$ with smooth, possibly nonclosed lower bound $-\theta$. Its tangent-bundle curvature bound is dominated by a fixed multiple of $\omega$. The averaging operators have a constant $K_0>0$ such that
\[
 R_tu+K_0t^2\downarrow u.
\]
In local base coordinates $z$ and a complex radius variable $w$, put $t=|w|$. Write $\zeta\in T_z^{1,0}X$ and $\eta\in\mathbb C$ for the base and radius components of a tangent vector, and set $\lambda(z,t)=t\partial_t(R_tu(z)+K_0t^2)$. With fixed constants $A,C>0$ and an error function $\epsilon_0(t)\geq0$ tending to zero, the theorem bounds the full base--fiber complex Hessian by the base form $-\theta-A\lambda(z,t)\omega-\epsilon_0(t)\omega$, together with mixed and vertical errors bounded by
$C(|\zeta||\eta|+|\eta|^2)$. Here $\lambda$ is nonnegative and increasing in $t$, and $\epsilon_0(t)\downarrow0$ as $t\downarrow0$. This is the classical regularization estimate for a smooth lower bound.

\Needspace{13\baselineskip}
\begin{lemma}\label{lem:general-kiselman-modulus}
For $K$ sufficiently large, there is an increasing function $\epsilon(\delta)\geq0$ tending to zero at zero such that, for every small $\delta>0$ and $c>0$,
\[
 \Phi_{c,\delta}
   =\inf_{0<t\leq\delta}
      \{R_tu+K(t^2-\delta^2)+K(t-\delta)-c\log(t/\delta)\}
\]
is bounded, quasi-plurisubharmonic, and satisfies
\[
 \theta+\ddc\Phi_{c,\delta}
                       \geq-(Ac+\epsilon(\delta))\omega.
\]
The infimum is attained away from zero. One may also use $K$ in the monotone smoothing inequality above.
\end{lemma}
\begin{proof}
Take $K\geq K_0$ and put $H(z,t)=R_tu(z)+Kt^2+Kt$. The function $H$ is increasing and strictly convex as a function of $\log t$, because $R_tu+K_0t^2$ has nonnegative increasing logarithmic derivative and the added terms have positive second logarithmic derivative. Its logarithmic derivative $\widetilde\lambda$ is at least $\lambda$.

The complex Hessian of $K|w|$ has a positive vertical term proportional to $K/|w|$. For sufficiently small $|w|=t$, it absorbs the vertical error and the mixed term, using
\[
 C|\zeta||\eta|
       \leq C' t|\zeta|^2+\frac{K}{8t}|\eta|^2.
\]
After increasing constants and replacing $\lambda$ by the larger $\widetilde\lambda$, the full base--fiber Hessian of $H$ is therefore bounded below by the base form
\[
 -\theta-(A\widetilde\lambda+\epsilon(t))\omega,
\]
where an increasing majorant of $\epsilon_0(t)+C't$ may be used.

Boundedness of $u$ makes the logarithmic penalty tend uniformly to $+\infty$ at zero. Thus the minimum occurs on a compact radial interval away from zero. Strict logarithmic convexity makes the minimizing radius unique. At an interior minimum $\widetilde\lambda=c$; at the upper cap $t=\delta$, one has $\widetilde\lambda\leq c$. No differentiability of the minimizing radius is needed.

For completeness, localize near any base point and fix $a>0$. Continuity of compact minimizers and of $\widetilde\lambda$ permits a radial interval containing all nearby minimizers on which $\widetilde\lambda\leq c+a$. In the cap case keep the original upper radius $\delta$; an open interval approaching the cap has the same infimum. The lower endpoint can be chosen below all minimizers by the uniform logarithmic penalty. Monotonicity of $\widetilde\lambda$ ensures its bound throughout the chosen interval. The minimum principle for rotation-invariant plurisubharmonic functions, applied after a local quadratic correction, now gives the lower curvature bound with $A(c+a)+\epsilon(\delta)$ and an arbitrarily small error from freezing the smooth base forms. Let these errors and $a$ tend to zero. This proves the asserted current inequality and quasi-plurisubharmonicity. Finally, $K\geq K_0$ preserves monotonicity and the pointwise limit of $R_tu+Kt^2$.
\end{proof}

The classical theorem supplies a modulus $\epsilon$, not a rate claimed here to be linear in $\delta$. The cancellation and kernel argument of \cref{lem:chern-smoothing-integral} still give
\[
 \|R_tu-u\|_{L^1(\mathrm{d}V)}\leq Ct^2,
 \qquad \|\nabla R_tu\|_\infty\leq C/t.
\]
Indeed the averaging map is unchanged, $u$ is bounded, and the monotone lower bound just proved is all that argument uses.


\subsection{A linear error for a smooth background}

Let $\theta$ be a smooth big real $(1,1)$-form. Choose a strict potential $\rho$ with analytic poles and smooth resolved remainder as in \cref{thm:resolvable-regularization-input}, and put $\Omega=X\setminus\{\rho=-\infty\}$. The next estimate refines the curvature error for normalized potentials locally bounded on $\Omega$.

\Needspace{24\baselineskip}
\begin{lemma}\label{lem:smooth-background-regularization-rate}
Let $a$ be quasi-plurisubharmonic on $X$, with $\sup_Xa=0$ and
$\theta+\ddc a\geq0$, and suppose that $a$ is locally bounded on $\Omega$.
There are constants $A,K,C>0$ such that the averaging operators $R_t$
from \cref{lem:chern-kiselman-input} satisfy
\[
 R_ta+Kt^2\downarrow a,\qquad
 \|R_ta-a\|_{L^1(\mathrm{d}V)}\leq Ct^2.
\]
For small $\delta>0$ and $c>0$, the transform on $\Omega$
\[
 \Phi_{c,\delta}
 =\inf_{0<t\leq\delta}
 \{R_ta+K(t^2-\delta^2)+K(t-\delta)-c\log(t/\delta)\}
\]
extends canonically to $X$ and satisfies
\begin{equation}\label{eq:smooth-background-transform-rate}
 a-K(\delta^2+\delta)\leq\Phi_{c,\delta}\leq R_\delta a\leq0,
 \qquad \theta+\ddc\Phi_{c,\delta}\geq-(Ac+C\delta)\omega.
\end{equation}
At each point of $\Omega$ the infimum is attained at a positive radius.
On every compact $K_1\Subset\Omega$, for sufficiently small $t$,
$\|\nabla R_ta\|_{L^\infty(K_1)}\leq C_{K_1}/t$.
\end{lemma}

\begin{proof}
Apply the classical Chern-connection regularization expansion to the closed current $\ddc a$ and the smooth lower bound $-\theta$. It gives monotonicity and logarithmic
convexity: for a constant $K_0>0$, set
$\lambda(x,t)=t\partial_t(R_ta(x)+K_0t^2)$. This function is nonnegative and increasing in $t$.
We first justify a linear error for this smooth lower bound, rather than
infer a rate from the unspecified modulus in that theorem.

Choose the auxiliary curvature form to be $A\omega$, with $A$ large
enough that the Chern curvature plus $A\omega\otimes\operatorname{Id}$
is nonnegative on every decomposable tensor. Equivalently, after
absorbing the normalization constants, writing $R_{j\bar k\ell\bar m}$ for the Chern curvature coefficients and $\xi,\zeta$ for tangent vectors in a unitary frame, the coefficients satisfy
\[
 \sum R_{j\bar k\ell\bar m}\xi_j\bar\xi_k\zeta_\ell\bar\zeta_m
 \geq-A|\xi|^2|\zeta|^2.
\]
This choice is available by compactness. Pairing this inequality with the nonnegative eigenvalues of a positive complex Hessian bounds its curvature contraction by $-A$ times its trace, with no further error term.

Fix $x\in X$ and a complex radius variable $w\ne0$, and put $t=|w|$. Write $\xi\in T_x^{1,0}X$ and $\tau\in\mathbb C$ for a base--radius tangent vector. In a uniformly sized coordinate chart about $x$,
write $\beta=\ddc|z|^2$ and
choose a quadratic $q_{x,t}$ with
\[
 \ddc q_{x,t}=\theta(x)+Bt\beta,
\]
where $B$ is fixed and large enough that $a+q_{x,t}$ is plurisubharmonic
on a ball containing the averaging support and its nearby base points.
This follows from the smooth variation bound
$\|\theta(y)-\theta(x)\|\leq C|y-x|$ on that ball of radius $O(t)$.
For the Hessian calculation at $(x,w)$ the quadratic is held fixed.
Its coefficients and all needed derivatives are uniformly bounded.

We estimate the errors in the classical regularization expansion.
Normalized quasi-plurisubharmonic compactness bounds $\|a\|_1$.
On each fixed larger chart, add one fixed strictly plurisubharmonic
quadratic $Q$ to $a$. If
$\nu=\ddc(a+Q)\wedge\beta^{n-1}$, the local Lelong mass inequality and
a fixed larger-ball mass bound give
$\nu(B_t(x))\leq Ct^{2n-2}$. The coefficients of
$\partial\bar\partial(a+q_{x,t})$ are controlled in total variation
by those of $\ddc(a+Q)$ plus a uniformly bounded smooth form. Write $|\partial\bar\partial(a+q_{x,t})|$ for the sum of the total variations of its coefficient measures. The rescaled averaging kernels are supported in $B(x,Ct)$ and bounded by a constant times $t^{-2n}$. Their integrals against this measure are therefore bounded by
\[
 C t^{-2n}|\partial\bar\partial(a+q_{x,t})|(B(x,Ct))
 \leq C't^{-2}.
\]
Thus the terms with a coefficient $O(t^3)$ contribute
$O(t)$; the order-two expansion has an $O(t)$ remainder as well.
The mixed and vertical contributions have bounds
$C(|\xi||\tau|+|\tau|^2)$. The curvature term is bounded below by
$-A\lambda(x,t)|\xi|^2$, using the preceding curvature choice.
The radial term of the same expansion identifies this mass with $\lambda$ up to $O(t^2)$.
Subtracting the regularization of $q_{x,t}$ restores $-\theta(x)$;
its base Hessian differs from $\ddc q_{x,t}$ by $O(t)$, its mixed and
vertical Hessians are bounded, and its radial logarithmic derivative
is $O(t^2)$. The extra $Bt\beta$ also costs only $O(t)$.
For a real $(1,1)$-form $\alpha$, write $\alpha[v]^2$ for its associated Hermitian quadratic form evaluated on $v$. Consequently, for a fixed constant $C_1>0$, the full base--fiber Hessian satisfies
\[
 \ddc_{x,w}R_{|w|}a[\xi,\tau]^2
 \geq-\theta_x[\xi]^2-(A\lambda(x,t)+Ct)|\xi|^2
       -C_1(|\xi||\tau|+|\tau|^2).
\]
All constants are independent of $x,t,c,\delta$ and uniform in the
normalized family. In particular, setting $\xi=0$ in the same local
expansion gives a uniform lower bound for the fiber Hessian. Adding a
fixed multiple of $|w|^2$ gives its subharmonicity, hence the
logarithmic convexity and monotonicity above with a uniform $K_0$.

Increase $K\geq K_0$ and add $K(t^2+t)$ to the regularization. The
positive vertical Hessian of $K|w|$ absorbs the mixed term by
\[
 C_1|\xi||\tau|\leq C't|\xi|^2+\frac{K}{8t}|\tau|^2,
\]
and absorbs the remaining bounded vertical error for small $t$.
The logarithmic derivative of the corrected function is at least
$\lambda$. The capped Kiselman minimum argument in
\cref{lem:general-kiselman-modulus} therefore gives the last inequality
in \eqref{eq:smooth-background-transform-rate}, now with error $C\delta$.
On compact subsets of $\Omega$, local boundedness of $a$ makes the
logarithmic penalty diverge uniformly at zero, so that the minimum
occurs away from zero. The same argument handles minimizers at the
upper cap without differentiating their radius.

The other bounds in \eqref{eq:smooth-background-transform-rate} follow
from monotonicity and evaluation at $t=\delta$. In particular the
transform is bounded above near the analytic complement. The removable
singularity theorem extends it, preserving the exact smooth lower
Hessian bound by local freezing and passage to the limit. Its lower
bound by $a-K(\delta^2+\delta)$ also extends by canonical representatives.

The adjoint computation in \cref{lem:chern-smoothing-integral} gives
$R_t^*1=1+O(t^2)$. It needs only $a\in L^1$, since
\[
 \left|\int_X(R_ta-a)\,\mathrm{d}V\right|
 \leq Ct^2\|a\|_1.
\]
Together with $R_ta-a\geq-Kt^2$, this proves the stated $L^1$ estimate.
Finally the smoothing kernel near $K_1$ samples a fixed larger compact
subset of $\Omega$, where $a$ is bounded. The same kernel derivative
estimate as in \cref{lem:chern-smoothing-integral} gives $C_{K_1}/t$.
\end{proof}
